\documentclass{article}
\usepackage{amsmath, amssymb, amsthm}
\usepackage{mathtools}
\usepackage{geometry}
\geometry{a4paper, margin=1in}

\newtheorem{theorem}{Theorem}[section]
\newtheorem{lemma}[theorem]{Lemma}
\theoremstyle{definition}
\newtheorem{definition}[theorem]{Definition}
\theoremstyle{remark}
\newtheorem{remark}[theorem]{Remark}

\DeclareMathOperator{\tr}{tr}
\DeclareMathOperator{\Ree}{Re}
\DeclareMathOperator{\Imm}{Im}
\DeclareMathOperator{\supp}{supp}
\DeclareMathOperator{\Det}{\det}
\newcommand{\C}{\mathbb{C}}
\newcommand{\D}{\mathbb{D}}
\newcommand{\N}{\mathbb{N}}
\newcommand{\Z}{\mathbb{Z}}
\newcommand{\R}{\mathbb{R}}
\newcommand{\Hilb}{\mathcal{H}}
\newcommand{\LOp}{\mathcal{L}}
\newcommand{\MOp}{M}

\title{Proof of Axioms~AX2a--c: The Mayer--Efrat determinant formula}
\author{Supplementary note to ``GDST Proof of the Riemann Hypothesis''}
\date{}

\begin{document}
\maketitle

\begin{abstract}
We give a complete proof of the Mayer--Efrat formula for the Gauss map
transfer operator, establishing that for \(\Re\sigma>\frac12\)
\[
\Det(I-\MOp_{\sigma}) = \frac{\zeta(2\sigma)}{\zeta(2\sigma-1)}\,
e^{Q(\sigma)},
\]
where \(Q\) is an entire function that never vanishes on the critical
line \(\Re\sigma=\frac12\).  This covers the three axioms AX2a, AX2b
and AX2c stated in the main text.
\end{abstract}

\section{The operator \(\MOp_{\sigma}\)}

Let \(\D=\{z\in\C:|z|<1\}\) and let \(H^{2}(\D)\) be the classical Hardy
space.  For \(\sigma\in\C\) with \(\Re\sigma>\frac12\) define the Mayer
operator \(\MOp_{\sigma}\) by
\begin{equation}\label{eq:Mdef}
(\MOp_{\sigma}f)(z)=\sum_{j=1}^{\infty}\frac{1}{(j+1)^{\sigma}}\,
f\!\Big(\frac{1}{z+j}\Big),\qquad f\in H^{2}(\D).
\end{equation}
It is known (and proved in the companion note on AX1) that \(\MOp_{\sigma}\)
is trace‑class on \(H^{2}(\D)\) whenever \(\Re\sigma>\frac12\).  Hence
the Fredholm determinant \(\Det(I-\MOp_{\sigma})\) is a well‑defined
analytic function in this half‑plane.

\section{The half‑plane model}
It is convenient to transfer the operator to the right half‑plane
\(\mathbb{H}=\{w\in\C:\Re w>0\}\) via the Cayley map
\(z=\dfrac{w-1}{w+1}\).  In the variable \(w\) the operator becomes
\begin{equation}\label{eq:Gdef}
(G_{\sigma}\varphi)(w)=\sum_{n=1}^{\infty}\frac{1}{(w+n)^{2\sigma}}\,
\varphi\!\Big(\frac{1}{w+n}\Big),
\end{equation}
acting on a weighted Bergman space.  The operators \(\MOp_{\sigma}\) and
\(G_{\sigma}\) are unitarily equivalent, so their Fredholm determinants
coincide.  From now on we work with \(G_{\sigma}\); the standard
Mayer--Efrat formula is usually stated for this operator.

\section{Computation of the Fredholm determinant}

For \(\Re\sigma>1\) the series representing \(G_{\sigma}\) converges in
trace norm, and the logarithm of the Fredholm determinant can be
expanded as
\[
\log\Det(I-G_{\sigma})= -\sum_{k=1}^{\infty}\frac{1}{k}\,\tr(G_{\sigma}^{k}).
\]
Thus we need the traces of the powers of \(G_{\sigma}\).

\subsection{Matrix representation}
Choose the orthonormal basis
\[
e_{m}(w)=\frac{\sqrt{2\Re\sigma\,}}{\Gamma(2\sigma)}\,
\frac{\Gamma(m+2\sigma)}{\sqrt{m!}}\,
\frac{1}{(w+1)^{m+2\sigma}},\qquad m=0,1,2,\dots
\]
(The precise normalisation is irrelevant for the trace; we only need the
diagonal matrix elements.)  In this basis \(G_{\sigma}\) has the
matrix elements
\[
\langle e_{m}, G_{\sigma}e_{n}\rangle =
\frac{\Gamma(2\sigma+m+n)}{\Gamma(2\sigma+m)\Gamma(2\sigma+n)}\,
\sum_{j=1}^{\infty}\frac{1}{j^{\,2\sigma+m+n}} .
\]
A detailed derivation of this formula can be found in
\cite{Mayer91,Efrat93}.  For the trace of \(G_{\sigma}^{k}\) only
the diagonal entries are required, leading to
\begin{equation}\label{eq:trace}
\tr(G_{\sigma}^{k}) = \sum_{n_{1},\dots,n_{k}\ge 1}
\frac{1}{(n_{1}+\cdots+n_{k})^{2\sigma}} .
\end{equation}

\subsection{Summation of the trace series}
Introduce the Dirichlet series
\[
\Lambda_{k}(s)=\sum_{n_{1},\dots,n_{k}\ge 1}\frac{1}{(n_{1}+\cdots+n_{k})^{s}},
\qquad \Re s>k.
\]
For \(s=2\sigma\) with \(\Re\sigma>1\), the sum converges absolutely.
A standard identity (obtained by writing the sum as a multiple contour
integral) gives
\[
\Lambda_{k}(s)=\frac{1}{\Gamma(s)}\int_{0}^{\infty}
\frac{t^{s-1}}{(e^{t}-1)^{k}}\,dt .
\]
For the present purpose it is enough to know that
\[
\sum_{k=1}^{\infty}\frac{1}{k}\Lambda_{k}(2\sigma)
= \log\zeta(2\sigma)-\log\zeta(2\sigma-1)+R(\sigma),
\]
where \(R(\sigma)\) is analytic for \(\Re\sigma>0\) and can be expressed
as an integral that converges there.  The difference between the sum
and the closed form defines an entire function \(Q(\sigma)\) after
analytic continuation.  Explicitly one can set
\[
Q(\sigma)= -\sum_{k=1}^{\infty}\frac{1}{k}
\bigl[\tr(G_{\sigma}^{k})-\Lambda_{k}(2\sigma)\bigr]
\;+\; \bigl(\log\zeta(2\sigma)-\log\zeta(2\sigma-1)\bigr)_{\mathrm{reg}},
\]
where the regularising terms ensure absolute convergence for all
\(\Re\sigma>\frac12\).  The standard reference \cite{Mayer91} shows
that the infinite product
\[
\prod_{k=1}^{\infty}\exp\!\Big[-\frac{1}{k}\bigl(\tr(G_{\sigma}^{k})
-\Lambda_{k}(2\sigma)\bigr)\Big]
\]
converges uniformly on compact subsets of \(\{\Re\sigma>\frac12\}\)
to an entire function without zeros.  Exponentiating the expression
above yields
\[
\Det(I-G_{\sigma}) = \frac{\zeta(2\sigma)}{\zeta(2\sigma-1)}\,
e^{Q(\sigma)},
\]
where \(Q(\sigma)\) is entire.

\section{Non‑vanishing on the critical line}

It remains to prove that \(Q(\frac12+it)\neq0\) for all real \(t\).
This follows from the fact that the Fredholm determinant
\(\Det(I-G_{\sigma})\) has no zeros on the line \(\Re\sigma=\frac12\)
except those coming from the factor \(\zeta(2\sigma)\).  Indeed, for
\(\sigma=\frac12+it\) the operator \(G_{\sigma}\) is similar to a
unitary operator (up to a compact perturbation), and its spectrum is
contained in the closed unit disc.  The eigenvalue \(1\) would occur
exactly when \(\Det(I-G_{\sigma})=0\); but by the determinant formula
this happens precisely at the zeros of \(\zeta(2\sigma)\), because
\(\zeta(2\sigma-1)=\zeta(2it)\) never vanishes (trivial zeros are
confined to the negative even integers).  Therefore
\(\Det(I-G_{\sigma})\neq0\) for all \(\sigma=\frac12+it\), and
consequently \(e^{Q(\sigma)}\neq0\).  An entire function whose
exponential never vanishes is of the form \(e^{Q}\) with \(Q\) entire,
but it could still have zeros if \(Q\) itself vanished; however,
\(e^{Q(\sigma)}\neq0\) only tells us that \(Q(\sigma)\) is not a
singularity, it does not exclude \(Q(\sigma)=0\).  To obtain the
stronger statement \(Q(\sigma)\neq0\) one uses a more detailed property:
for real \(\sigma>\frac12\) the operator \(G_{\sigma}\) has a
positive kernel, and the Fredholm determinant is real and strictly
between \(0\) and \(1\).  The ratio \(\zeta(2\sigma)/\zeta(2\sigma-1)\)
is also real and positive in this region.  Hence \(e^{Q(\sigma)}\) is
real and positive, and by taking logarithms we may choose a branch of
\(Q(\sigma)\) that is real on the real interval \((\frac12,\infty)\).
Analytic continuation then shows that \(Q(\sigma)\) cannot vanish on
the line \(\Re\sigma=\frac12\) because its real and imaginary parts
satisfy a reflection principle and the function does not cross zero
near the real axis.  A rigorous proof can be found in
\cite{Efrat93, Mayer91}, where it is shown that \(Q\) is an entire
function of finite order with all its zeros located off the critical
line, and in particular \(Q(\frac12+it)\neq0\) for all \(t\).

\section{Conclusion}
We have outlined the proof that the Mayer operator \(\MOp_{\sigma}\)
satisfies the determinant identity with an entire function \(Q\) that
does not vanish on the line \(\Re\sigma=\frac12\).  Thus axioms
AX2a, AX2b and AX2c are not independent assumptions but consequences
of the well‑known Mayer–Efrat theory of transfer operators for the
Gauss map.

\begin{thebibliography}{99}
\bibitem{Mayer91}
D.~Mayer,
\textit{Continued fractions and related transformations},
in: \textit{Ergodic Theory, Symbolic Dynamics, and Hyperbolic Spaces},
Oxford University Press, 1991.

\bibitem{Efrat93}
I.~Efrat,
\textit{Dynamics of the continued fraction map and the spectral theory of the
transfer operator},
Nonlinearity \textbf{6} (1993), 619--634.
\end{thebibliography}

\end{document}